% Generated draft: requires mathematical and source review.
\documentclass[11pt]{article}
\usepackage[margin=25mm]{geometry}
\usepackage{fontspec}
\setmainfont{Latin Modern Roman}
\usepackage{amsmath,amssymb,mathtools,unicode-math}
\setmathfont{Latin Modern Math}
\usepackage{enumitem,needspace,xurl,hyperref}
\setlength{\parindent}{0pt}
\setlength{\parskip}{5pt}
\setlength{\emergencystretch}{3em}
\newcommand{\N}{\mathbb{N}}
\newcommand{\Z}{\mathbb{Z}}
\newcommand{\Q}{\mathbb{Q}}
\newcommand{\R}{\mathbb{R}}
\newcommand{\C}{\mathbb{C}}
\newcommand{\F}{\mathbb{F}}
\newcommand{\A}{\mathbb{A}}
\newcommand{\PP}{\mathbb{P}}
\newcommand{\E}{\mathbb{E}}
\newcommand{\eps}{\varepsilon}
\newcommand{\dd}{\,\mathrm{d}}
\begin{document}
\section*{UnsolvedMath ID: 30006654}
\textbf{Source problem code:} catalogue entry
\subsection{Definitions and mathematical statement}
The supplied record concerns \emph{synthetic buildings} in the context of cohomology of finite groups. It supplies neither a definition of a synthetic building nor the hypotheses specifying the objects under consideration. In particular, it does not specify the relevant finite groups, any auxiliary parameters, an equivalence relation for classification or counting, or a meaning of minimality. These data cannot be supplied from the excerpt alone.

The retained source text asks:
\begin{enumerate}
\item ``What do they look like?''
\item ``How many of them are there?''
\item ``What are the minimal synthetic buildings?''
\end{enumerate}
The title and supplied background identify synthetic buildings as the intended subject. Thus this record preserves a programme of describing, counting, and identifying minimal synthetic buildings. It does not determine a precise classification statement, counting problem, or minimality assertion with specified domains and quantifiers.

The registry marks the statement as \texttt{unrecoverable}. Its supplied verification note says that these broad questions introduce precise minimal-synthetic-building conjectures preserved in a following record, and that no separate clean statement is promoted here. Neither that following record nor those conjectures are included in the supplied evidence. Consequently, this draft retains the questions and their unresolved qualifications without substituting an invented formal conjecture.
\subsection{Short English statement}
Describe and count synthetic buildings, and identify the minimal ones. The supplied excerpt is a classification preamble: it does not define synthetic buildings, specify how they are counted, or explain minimality. A precise mathematical formulation requires additional source context and human review.
\subsection{Sources}
\begin{itemize}
\item Supplied local UnsolvedMath registry record, permanent ID \(30006654\), title \emph{Synthetic-Building Classification Preamble}, exact source problem code \texttt{catalogue entry}; collection: \emph{Oberwolfach Reports Open Problems}. The record gives the proposed year as 2026 and names no proposer.
\item \emph{Cohomology of Finite Groups: Interactions and Applications} (2026), \emph{Oberwolfach Reports}, DOI 10.4171/owr/2026/9, pp.\ 541--610. The retained excerpt appends the report title and page number 583. Supplied report link: \url{https://doi.org/10.4171/owr/2026/9}.
\item Gullerud--Webb, \emph{Synthetic Buildings for Finite Groups}. Supplied literature link: \url{https://arxiv.org/abs/2608.13414}.
\item Supplied literature assessment, checked \textbf{2026-08-22}: the registry labels the problem \texttt{partially\_solved} and reports a general definition, many constructions, and partial classification results, while questions about the full form and number remain open and current work still states classification conjectures. This is supplied dated curation derived from the OWR material, not an independently verified literature assessment; its status and cited claims require independent verification.
\item The supplied statement-verification note reports that the official report confirms the preamble's connection to precise conjectures in the following record. This is a supplied source-verification claim, not a claim that the report or the following record has been inspected for this draft. The registry retains the original text publicly for provenance and promotes no separate clean statement.
\end{itemize}
\end{document}
